% Source-aligned English for printed p.203.
\clearpage
\SourcePage{203}

\LayerTarget{tohoku:remark:5.2.4:2}
\SourceSideDiagram{%
\LayerTarget{tohoku:diagram:five-term-comparison-5.2}
\setlength{\unitlength}{0.46mm}
\begin{picture}(105,78)
\put(43.5,75){\makebox(0,0){$0$}}
\put(43.5,71){\vector(0,-1){6}}
\put(43.5,61){\makebox(0,0){$A'$}}
\put(43.5,57){\vector(0,-1){9}}
\put(40.5,52){\makebox(0,0)[r]{$\alpha$}}
\put(1.5,44){\makebox(0,0){$0$}}
\put(3.75,44){\vector(1,0){9}}
\put(16.5,44){\makebox(0,0){$A$}}
\put(20.25,44){\vector(1,0){18.75}}
\put(30,48){\makebox(0,0){$\alpha'$}}
\put(43.5,44){\makebox(0,0){$H^1$}}
\put(48,44){\vector(1,0){15.75}}
\put(56.25,39){\makebox(0,0){$\beta'$}}
\put(68.25,44){\makebox(0,0){$B$}}
\put(72,44){\vector(1,0){9}}
\put(76.5,39){\makebox(0,0){$\gamma'$}}
\put(84.75,44){\makebox(0,0){$C$}}
\put(88.5,44){\vector(1,0){8.25}}
\put(92.25,39){\makebox(0,0){$\delta'$}}
\put(101.25,44){\makebox(0,0){$D$}}
\put(43.5,40){\vector(0,-1){7}}
\put(46.5,36){\makebox(0,0)[l]{$\beta$}}
\put(43.5,29){\makebox(0,0){$B'$}}
\put(43.5,25){\vector(0,-1){5}}
\put(46.5,22){\makebox(0,0)[l]{$\gamma$}}
\put(43.5,17){\makebox(0,0){$C'$}}
\put(43.5,13){\vector(0,-1){5}}
\put(46.5,10){\makebox(0,0)[l]{$\delta$}}
\put(43.5,5){\makebox(0,0){$D'$}}
\put(46.5,59){\vector(3,-2){18.75}}
\put(57,58){\makebox(0,0){$u$}}
\put(19.5,42){\vector(3,-2){21}}
\put(29.25,31){\makebox(0,0){$u'$}}
\end{picture}%
}
\noindent and that of $G$. More generally, suppose that two five-term exact
sequences are given, as in the present case (see the diagram below).
\LayerTarget{tohoku:construction:five-term-comparison-maps}
Set $u=\beta'\alpha$ and $u'=\beta\alpha'$. Then
$A\cap A'=\operatorname{Ker}u=\operatorname{Ker}u'$. There are natural
homomorphisms
\[
\begin{aligned}
\beta\beta'^{-1}&:\operatorname{Ker}\gamma'\to
  \operatorname{Coker}u',\\
\beta'\beta^{-1}&:\operatorname{Ker}\gamma\to
  \operatorname{Coker}u.
\end{aligned}
\]
These homomorphisms give rise to the following two exact sequences (in which
$H^1$ no longer appears):

\LayerTarget{tohoku:equation:5.2.12}
\[
\tag{5.2.12}
\left\{
\begin{aligned}
0\to{}&\operatorname{Ker}u\to A'
\to\operatorname{Ker}\gamma'
\to\operatorname{Coker}u'\to C'\to D',\\[-0.2em]
&\hphantom{\operatorname{K}}\mathrel{\Vert}\\[-0.45em]
0\to{}&\operatorname{Ker}u'\to A
\to\operatorname{Ker}\gamma
\to\operatorname{Coker}u\to C\to D.
\end{aligned}
\right.
\]
\LayerTarget{tohoku:example:five-term-comparison}
For example, let $u$ be the natural homomorphism from $H^1(Y,A^G)$ to
$H^1(X,A)^G$, and let $u'$ be the natural homomorphism from
$H^1(G,\Gamma(A))$ to $H^0(Y,\mathbf H^1(G,A))$. Then \textit{the kernel of
$u$ is isomorphic to that of $u'$, whereas the image of $u$ is isomorphic to
the kernel of the natural homomorphism from
$\operatorname{Ker}(H^1(X,A)^G\to H^2(G,\Gamma(A)))$ to
$\operatorname{Coker}u'$.}

\LayerTarget{tohoku:remark:5.2.4:3}
3. The hypotheses of Proposition 5.2.2 hold in most cases that
arise naturally. Its interest lies chiefly in formula (5.2.11), which, when
$G$ is finite, can also be written as

\LayerTarget{tohoku:equation:5.2.11-bis}
\[
\mathbf H^q(G,A)(y)=H^q(G,f_*(A)(y))\qquad y\in Y
\tag{5.2.11 bis}
\]
\LayerTarget{tohoku:application:finite-group-algebraic-variety}
This gives an explicit computation of the sheaves $\mathbf H^q(G,A)$.
Proposition 5.2.2 applies whenever $G$ is finite and $X$ is separated. Another
important case is that in which $G$ is a finite group of automorphisms of an
abstract algebraic variety $X$ and $Y=X/G$ is an algebraic variety (for
example, $G$ is the Galois group of a covering $X$, ramified or not, of a
normal variety $Y$). The hypotheses of Proposition 5.2.2 then hold if $A$ is
a \textit{coherent algebraic $G$-sheaf}, or if $A$ is the multiplicative
sheaf $\mathbf O_X^*$ of germs of invertible regular functions on $X$.
Moreover, if $X$ is unramified over $Y$, then $\mathbf H^n(G,A)=0$ for $n>0$
when $A$ is a coherent algebraic G-sheaf, which allows Proposition 5.2.3 to
be applied. In the same case one can also show that
$\mathbf H^1(G,\mathbf O_X^*)=0$. Arithmetic varieties could also be
considered; the same results hold.

\LayerTarget{tohoku:section:5.3}
\bigskip
\textbf{5.3. Case of a discontinuous group of homeomorphisms.} For
brevity, and by an abuse of language, we shall say that $G$ is a
\textit{discontinuous group of homeomorphisms of $X$} if it satisfies the
following condition:

\LayerTarget{tohoku:condition:5.3:D}
\textit{\textup{(D)} For every $x\in X$, the stabilizer $G_x$ of $x$ is
finite, and there exists a neighborhood $V_x$ of $x$ such that, for every
$g\in G$ not in $G_x$, one has $g\mathbin{.}V_x\cap V_x=\phi$.}

\LayerTarget{tohoku:remark:condition-D-basic-properties:p1}
We may then clearly suppose that $V_x$ is open and, moreover, that
$g\mathbin{.}V_x=V_x$ for $g\in G_x$ (replacing $V_x$, if necessary, by the
intersection of the $g\mathbin{.}V_x$ for $g\in G_x$). Condition (D) holds, for
example, if $G$ is finite and $X$ is separated, as is immediately clear. If
$X$ is an irreducible algebraic variety endowed with its

\TohokuContinuation{204}
